\documentclass{article}
\usepackage[T1]{fontenc}
\usepackage{lmodern}
\usepackage{amsmath,amssymb,amsthm,mathtools}
\usepackage[hidelinks]{hyperref}
\newtheorem*{maintheorem}{Main theorem}
\DeclareMathOperator{\wt}{wt}
\DeclareMathOperator{\rowsp}{row}
\DeclareMathOperator{\supp}{supp}
\title{Exact CSS parameters under uniform qubit replication}
\author{}
\date{}
\begin{document}
\maketitle

\textbf{Grade: THEOREM.} Complete hand proof, three independent hostile reviews,
adversarial author rereading, and guarded Lean replay.

\begin{maintheorem}\label{thm:main}
Let a binary CSS code have a finite physical set \(Q\), X-check matrix \(A\),
and Z-check matrix \(B\), with \(AB^{\mathsf T}=0\). Let \(n=|Q|\), and let
\begin{equation}\label{eq:old-params}
r_X=\dim\rowsp(A),\qquad
r_Z=\dim\rowsp(B),\qquad k=n-r_X-r_Z.
\end{equation}

Fix a nonempty finite set \(I\), choose \(i_0\in I\), and write \(r=|I|\).
Replace each physical coordinate \(q\) by the \(r\) coordinates \((i,q)\).
The following displayed rows define the replicated code:
\begin{itemize}
\item For every old X row \(a\), put its entries on the anchor coordinates
  \((i_0,q)\), with zero entries on every other replica.
\item For every \(q\in Q\) and \(i\in I\setminus\{i_0\}\), add the X row
  \(e_{(i,q)}+e_{(i_0,q)}\).
\item For every old Z row \(b\), put the entry \(b(q)\) at every \((i,q)\).
\end{itemize}

Then these rows commute. Their ranks and encoded dimension are
\begin{equation}\label{eq:new-params}
r'_X=r_X+n(r-1),\qquad r'_Z=r_Z,\qquad k'=k.
\end{equation}

If both original sector distances exist, then both replicated sector
distances exist and satisfy
\begin{equation}\label{eq:main-distances}
\boxed{d'_X=d_X,\qquad d'_Z=r d_Z.}
\end{equation}

If the old X/Z check weights are bounded by \(w_X,w_Z\), respectively, the
new bounds are \(\max(w_X,2)\) and \(r w_Z\). If an old qubit belongs to at
most \(\delta_X\) X checks and \(\delta_Z\) Z checks, every new qubit belongs
to at most \(\delta_X+r\) X checks and \(\delta_Z\) Z checks. Thus for each
fixed positive \(r\), this operation takes a qLDPC family to a qLDPC family,
preserves its encoded dimension, and gives its exact asymmetric distances.
The theorem does not claim uniform sparsity if \(r\) itself grows.

The statements allow dependent original rows. Neither the rank result nor
the sparsity result assumes the existence of any nontrivial logical.
\end{maintheorem}

\section{Binary vector spaces and distance}\label{sec:binary}

Write \(F=\{0,1\}\) for the field with addition and multiplication modulo
two. In particular \(1+1=0\). A vector on a finite set \(Q\) is a function
\begin{equation}\label{eq:vector}
x:Q\longrightarrow F.
\end{equation}

Its support is \(\{q:x(q)\ne0\}\), and its Hamming weight \(\wt(x)\)
is the number of elements in this support. The vector \(e_q\) has value one
at \(q\) and zero at every other coordinate. The dot product is
\begin{equation}\label{eq:dot}
x\cdot z=\sum_{q\in Q}x(q)z(q).
\end{equation}

The row space of a matrix is the space of binary linear combinations of
its rows. Write
\begin{equation}\label{eq:spaces}
U=\rowsp(A),\quad V=\rowsp(B),\quad
S_X=\ker B,\quad S_Z=\ker A.
\end{equation}

The equation \(AB^{\mathsf T}=0\) says that every row of \(A\) has dot
product zero with every row of \(B\). Expanding linear combinations of
rows therefore gives \(U\subseteq S_X\) and \(V\subseteq S_Z\).

A nontrivial X logical is a vector in \(S_X\setminus U\), and a nontrivial
Z logical is a vector in \(S_Z\setminus V\). If the relevant set is
nonempty, its distance is the smallest Hamming weight in that set.
These sets are finite, so a smallest weight is attained. We assign no
distance to an empty set. A distance assertion in this document includes
both a vector attaining the stated value and a proof that no nontrivial
logical has smaller weight.

\section{Three maps and their identities}\label{sec:maps}

For \(y\in F^{I\times Q}\) and \(x,z\in F^Q\), define
\begin{equation}\label{eq:maps}
(Py)(q)=\sum_{i\in I}y(i,q),\qquad
(Rx)(i,q)=x(q),\qquad
(Ex)(i,q)=
\begin{cases}
x(q),&i=i_0,\\
0,&i\ne i_0.
\end{cases}
\end{equation}

Each is a linear map because each output coordinate is a fixed sum of
input coordinates. For every \(x\), the sum defining \(PEx\) has one
possibly nonzero term, at \(i_0\), and equals \(x\). Thus \(PE\) is the
identity and \(P\) is surjective. If \(Rx=Rz\), evaluating at \((i_0,q)\)
for every \(q\) gives \(x=z\); hence \(R\) is injective.

Expansion of the two finite sums gives the dot-product identity
\begin{equation}\label{eq:adjoint}
y\cdot Rz
=\sum_{q\in Q}\sum_{i\in I}y(i,q)z(q)
=\sum_{q\in Q}(Py)(q)z(q)
=Py\cdot z.
\end{equation}

The three weight relations are
\begin{equation}\label{eq:weights}
\wt(Ex)=\wt(x),\qquad
\wt(Rx)=r\wt(x),\qquad
\wt(Py)\le\wt(y).
\end{equation}

For the first equality, \(q\mapsto(i_0,q)\) bijects the two supports.
For the second equality, the support of \(Rx\) is exactly
\begin{equation}\label{eq:repeat-support}
I\times\supp(x).
\end{equation}

For the inequality, if \((Py)(q)\ne0\), at least one term \(y(i,q)\) is
nonzero; a sum all of whose terms are zero is zero. Consequently
\begin{equation}\label{eq:parity-support}
\supp(Py)
\subseteq\{q:\text{there exists }i\text{ with }y(i,q)\ne0\}.
\end{equation}

The set on the right is the image of \(\supp(y)\) under
projection to \(Q\). An image of a finite set has no more elements than
that set, proving the inequality. This argument holds for even and odd
\(r\); it does not use \(PR\) being the identity.

\section{Exact stabilizer spaces}\label{sec:stabilizers}

Let \(p_{i,q}=e_{(i,q)}+e_{(i_0,q)}\). For \(i=i_0\), this is zero. We may
include these zero rows in a matrix without changing its row space,
commutation, or nonzero row incidences. This is the convention used in
the Lean matrix definition.

The parity of every \(p_{i,q}\) is zero. At base coordinate \(q\), its two
contributions add to zero; at each other base coordinate both are zero.
For any \(y\in F^{I\times Q}\), the following is an exact vector identity:
\begin{equation}\label{eq:star-sum}
\sum_{i\in I}\sum_{q\in Q}y(i,q)p_{i,q}=y+EPy.
\end{equation}

To check \eqref{eq:star-sum}, expand its left side using the definition of \(p_{i,q}\).
The terms \(y(i,q)e_{(i,q)}\) sum to \(y\), since exactly the term indexed
by a coordinate contributes at that coordinate. The remaining terms
are all on the anchor. Their coefficient at \((i_0,q)\) is
\begin{equation}\label{eq:anchor-coefficient}
\sum_i y(i,q)=(Py)(q),
\end{equation}

so they sum to \(EPy\). This proves \eqref{eq:star-sum} at every coordinate. In particular,
when \(Py=0\), the star rows span \(y\). Thus they span exactly \(\ker P\).

Write \(U',V'\) for the two replicated stabilizer spaces. We claim
\begin{equation}\label{eq:stabilizers}
U'=P^{-1}(U),\qquad V'=R(V).
\end{equation}

Each anchor X row has parity equal to its old row in \(U\); every star
row has parity zero in \(U\). Linearity proves \(U'\subseteq P^{-1}(U)\).
For the reverse inclusion, if \(Py\in U\), express \(Py\) as a sum of old
X rows. Applying \(E\) shows \(EPy\in U'\). Identity \eqref{eq:star-sum} shows
\begin{equation}\label{eq:star-member}
y+EPy\in U'.
\end{equation}

Adding these two vectors gives \(y\in U'\), because \(EPy+EPy=0\).
For the Z equality, each new Z row is the repeat of an old Z row.
Linearity of \(R\) implies that all their linear combinations are exactly
the repeats of the old linear combinations, which is \(R(V)\).

For \(u'\in U'\) and \(v'\in V'\), choose \(v\in V\) with \(v'=Rv\).
Then \(Pu'\in U\) by \eqref{eq:stabilizers}, and \eqref{eq:adjoint} gives
\begin{equation}\label{eq:commutation}
u'\cdot v'=Pu'\cdot v=0.
\end{equation}

Thus the displayed checks commute and define a CSS code.

\section{Exact cycle spaces}\label{sec:cycles}

An X vector \(y\) has zero dot product with every replicated Z row
precisely when \(Py\) has zero dot product with every old Z row, by
\eqref{eq:adjoint}.
Therefore its cycle space is
\begin{equation}\label{eq:x-cycles}
S'_X=P^{-1}(S_X).
\end{equation}

A Z vector \(z'\) has zero dot product with the star row \(p_{i,q}\)
precisely when
\begin{equation}\label{eq:star-syndrome}
z'(i,q)+z'(i_0,q)=0,
\end{equation}

equivalently \(z'(i,q)=z'(i_0,q)\). These equations for every \(i,q\)
force \(z'=Rz\), where \(z(q)=z'(i_0,q)\). Conversely every \(Rz\)
satisfies all of them. Its dot product with an anchored old row \(Ea\)
is \(a\cdot z\), by \eqref{eq:adjoint} and \(PEa=a\). Thus all anchored X syndromes
vanish precisely when \(z\in S_Z\). We have proved
\begin{equation}\label{eq:z-cycles}
S'_Z=R(S_Z).
\end{equation}

Equations \eqref{eq:stabilizers}, \eqref{eq:x-cycles}, and \eqref{eq:z-cycles} yield the
following
complete logical description:
\begin{equation}\label{eq:x-logicals}
y\in S'_X\setminus U'
\quad\Longleftrightarrow\quad Py\in S_X\setminus U,
\end{equation}
\begin{equation}\label{eq:z-logicals}
z'\in S'_Z\setminus V'
\quad\Longleftrightarrow\quad
z'=Rz\text{ for a unique }z\in S_Z\setminus V.
\end{equation}

For \eqref{eq:z-logicals}, the uniqueness and preservation of nonmembership follow from
injectivity of \(R\): if \(Rz=Rv\) with \(v\in V\), then \(z=v\in V\).
All directions of \eqref{eq:x-logicals} follow directly from the two inverse images.

\section{Exact distances}\label{sec:distances}

Suppose \(d_X\) exists. For every replicated nontrivial X logical \(y\),
equations \eqref{eq:x-logicals} and \eqref{eq:weights} give
\begin{equation}\label{eq:x-lower}
d_X\le\wt(Py)\le\wt(y).
\end{equation}

Choose an original X logical \(x\) of weight \(d_X\). The vector \(Ex\)
is a replicated nontrivial X logical because \(PEx=x\), and it has weight
\(d_X\) by \eqref{eq:weights}. Both inequalities necessary for \(d'_X=d_X\) are proved.

Suppose \(d_Z\) exists. By \eqref{eq:z-logicals} every replicated nontrivial Z logical has
the unique form \(Rz\) for an old nontrivial Z logical \(z\). Equation \eqref{eq:weights}
then gives
\begin{equation}\label{eq:z-lower}
\wt(Rz)=r\wt(z)\ge r d_Z.
\end{equation}

Repeating an old Z logical of weight \(d_Z\) attains \(r d_Z\). Thus
\(d'_Z=r d_Z\). These arguments also show that the relevant logical set
is empty in the new code exactly when it was empty in the old code.

\section{Ranks and encoded dimension}\label{sec:ranks}

Because \(PE\) is the identity, \(F^{I\times Q}\) decomposes as
\begin{equation}\label{eq:decomposition}
F^{I\times Q}=E(F^Q)\mathbin{\oplus}\ker P.
\end{equation}

Indeed \(y=EPy+(y+EPy)\), and applying \(P\) to the second summand gives
zero. If \(Ex\in\ker P\), then \(x=PEx=0\), proving that the intersection
is zero. A basis of \(F^Q\) consists of its \(n\) unit vectors, and a
basis of \(F^{I\times Q}\) consists of its \(rn\) unit vectors. An
injective linear map takes a basis to a basis of its image, and bases
of two subspaces with zero intersection concatenate to a basis of
their sum: spanning and linear independence follow by projecting a
linear relation to each direct summand. Consequently \eqref{eq:decomposition} implies
\begin{equation}\label{eq:kernel-dimension}
\dim\ker P=rn-n=n(r-1).
\end{equation}

The same decomposition restricted to \(P^{-1}(U)\) gives
\begin{equation}\label{eq:restricted-decomposition}
P^{-1}(U)=E(U)\mathbin{\oplus}\ker P.
\end{equation}

The proof is the decomposition above with \(Py\in U\); intersection
zero was already proved. Since \(E\) is injective, \eqref{eq:stabilizers} implies
\begin{equation}\label{eq:x-rank}
\dim U'=\dim U+\dim\ker P=r_X+n(r-1).
\end{equation}

Injectivity of \(R\) and \(V'=R(V)\) give \(\dim V'=\dim V=r_Z\).
Subtracting the two ranks from \(rn\) gives
\begin{equation}\label{eq:encoded-dimension}
k'=rn-\bigl(r_X+rn-n\bigr)-r_Z=n-r_X-r_Z=k.
\end{equation}

\section{Sparsity and families}\label{sec:sparsity}

An anchored old X row has its original weight by \eqref{eq:weights}. A star row has
weight two when \(i\ne i_0\): its two indicated coordinates are
different, so both entries are one, and all other entries are zero.
When \(i=i_0\), its weight is zero. These two cases are exhaustive.
Hence every X row has weight at most \(\max(w_X,2)\). A repeated Z row
has weight \(r\) times its old weight, so every Z row has weight at most
\(r w_Z\).

For a fixed new physical coordinate \((j,q)\), the old anchored X rows
contribute at most \(\delta_X\) incidences: only when \(j=i_0\) can any
contribute. Among star rows, a row \(p_{i,e}\) can contribute only when
\(e=q\), because both of its coordinates have base coordinate \(e\).
There are at most \(r\) choices of \(i\), so these rows contribute at
most \(r\) incidences. This proves the bound \(\delta_X+r\). A Z row is
nonzero at \((j,q)\) exactly when its original row is nonzero at \(q\),
giving exactly the original Z incidence count and the bound
\(\delta_Z\).

A qLDPC family here means a collection of CSS codes whose check row
weights and numbers of incident checks per qubit are bounded by
constants independent of the family index. For a fixed \(r\), the four
bounds just proved are constants whenever the original four bounds
are constants. The resulting physical size is \(rn\), encoded dimension
is \(k\), and sector distances are \(d_X,r d_Z\). This proves the stated
family conclusion without a numerical certificate or a search.

\section{Formal correspondence}\label{sec:formal}

\begin{sloppypar}
\texttt{lean/QldpcVerification/QubitReplication.lean} defines the maps as
\texttt{parity}, \texttt{replicate}, and \texttt{embed}; the displayed matrices as
\texttt{HX} and \texttt{HZ};
and proves \texttt{rowspace\_HX}, \texttt{rowspace\_HZ}, \texttt{kernel\_HX},
\texttt{kernel\_HZ}, \texttt{commutation}, \texttt{exact\_distances},
\texttt{ranks}, \texttt{encoded\_dimension}, \texttt{row\_weights},
\texttt{x\_column\_weight\_le}, and \texttt{z\_column\_weight}.
\end{sloppypar}

The formal exact-distance predicate is \texttt{Components.Distance}: it contains
both attainment and a universal lower bound. No generated certificate
or concrete code instance enters this proof.

\section{Attribution}\label{sec:attribution}

This is repetition-code concatenation, a known construction. Daniel Gottesman,
\emph{Stabilizer Codes and Quantum Error Correction} (1997), Sections 3.5 and 6.1,
provides the general stabilizer-concatenation construction
(\href{https://thesis.caltech.edu/2900/02/THESIS.pdf}{primary thesis}).
The proof above derives all required facts independently; no priority claim is made.

\end{document}
